% Generated draft: requires mathematical and source review.
\documentclass[11pt]{article}
\usepackage[margin=25mm]{geometry}
\usepackage{fontspec}
\setmainfont{Latin Modern Roman}
\usepackage{amsmath,amssymb,mathtools,unicode-math}
\setmathfont{Latin Modern Math}
\usepackage{enumitem,needspace,xurl,hyperref}
\setlength{\parindent}{0pt}
\setlength{\parskip}{5pt}
\setlength{\emergencystretch}{3em}
\newcommand{\N}{\mathbb{N}}
\newcommand{\Z}{\mathbb{Z}}
\newcommand{\Q}{\mathbb{Q}}
\newcommand{\R}{\mathbb{R}}
\newcommand{\C}{\mathbb{C}}
\newcommand{\F}{\mathbb{F}}
\newcommand{\A}{\mathbb{A}}
\newcommand{\PP}{\mathbb{P}}
\newcommand{\E}{\mathbb{E}}
\newcommand{\eps}{\varepsilon}
\newcommand{\dd}{\,\mathrm{d}}
\begin{document}
\section*{UnsolvedMath ID: 10400195}
\textbf{Source problem code:} catalogue entry
\subsection{Definitions and mathematical statement}
Consider pairs \((M,s)\), where \(M\) is a closed oriented three-manifold and \(s\) is a \(\operatorname{spin}^c\) structure, that is, a lift of its oriented frame bundle through \(\operatorname{Spin}^c(3)\to SO(3)\), with \(\operatorname{Spin}^c(3)=(\operatorname{Spin}(3)\times U(1))/\{(1,1),(-1,-1)\}\). Isomorphisms must carry the specified structure to the specified structure; the structure is part of the object.

Let \(\sim_{Y_d^c}\) be the spin-\(c\) clasper-equivalence relation intended in the source. Deloup and Massuyeau ask to describe the quotient set
\[\mathcal Q_d^c=\{(M,s):M\text{ a closed three-manifold},\ s\in\operatorname{Spin}^c(M)\}/\!\sim_{Y_d^c},\]
with particular emphasis on \(d=2\) and \(d=3\). A complete invariant description would specify a set \(A_d\), its realizable elements and an assignment \(I_d\) satisfying
\[(M,s)\sim_{Y_d^c}(M',s')\quad\Longleftrightarrow\quad I_d(M,s)=I_d(M',s'),\]
and identify its image, equivalently describe \(\mathcal Q_d^c\) up to a stated bijection. This clarifies completeness without requiring a finite list of invariants or an algorithm, neither of which the source demands. The task is about a quotient \emph{set}; no connected-sum group structure is imposed. No homology-sphere restriction appears.

The local moves, degree convention and transport of spin-\(c\) structures in \(Y_d^c\)-equivalence are not supplied, so the equivalence cannot be replaced by ordinary \(Y_d\)-equivalence on underlying manifolds. The permitted range of \(d\), connectedness convention and initial isomorphism category require surrounding definitions.
\subsection{Short English statement}
Classify closed three-manifolds together with their spin-\(c\) structures under the specified \(Y_d^c\)-equivalence, especially in degrees two and three. The structure-transport rules for that equivalence are absent.
\subsection{Sources}
\begin{itemize}
\item Local UnsolvedMath registry, permanent ID \(10400195\), exact source code \texttt{catalogue entry}.
\item F. Deloup and G. Massuyeau, attribution in Ohtsuki, \emph{Problems on invariants of knots and 3-manifolds} (2002), Problem 10.20, PDF page 154: \url{https://msp.org/gtm/2002/04/gtm-2002-04-024s.pdf}. Registry access 2026-07-29; pdf-font-structure extraction.
\end{itemize}
\end{document}
